\newpage
\section{Omitted Details in \pref{sec:pre}} \label{app: recovering}
In this section, we show that the algorithms in \cite{xu2023bayesian} and \cite{liu2025decision} are special cases of \pref{alg:general}. 

\subsection{Recovering \pref{thm: air phi}} The decision rule of \cite{liu2025decision}'s algorithm corresponds to \pref{eq:minmax} with $D^\pi(\nu\|\rho) = \E_{M\sim \nu}\E_{o\sim M(\cdot|\pi)}[\KL(\nu_{\bo{\phi}}(\cdot|\pi,o), \rho)]$. It can be shown that $\breg_{D^\pi(\cdot\|\rho)}(\nu, \nu')=\E_{M \sim \nu}\E_{o \sim M(\cdot|\pi)}\big[\KL(\nu_{\bo{\phi}}(\cdot|\pi, o), \nu'_{\bo{\phi}}(\cdot|\pi, o))\big]$ in this case. Furthermore, notice that when $\nu=\delta_{M_t, \picirc}$, we have $\nu_{\bo{\phi}}(\cdot|\pi, o)=\delta_{\phi^\star}$ according to \pref{def: restricted env}. Thus, the estimation error term in \pref{eq:minimax-guarantee} in \cite{liu2025decision}'s algorithm is 
\begin{align*}
    \E[\est] &= \E\left[\sum_{t=1}^T \left( \KL(\delta_{\phi^\star}, \rho_t) - \E_{o\sim M_t(\cdot|\pi_t)} \Big[\KL(\delta_{\phi^\star},  (\nu_t)_{\bo{\phi}}(\cdot|\pi_t,o))\Big]\right) \right] \\
    &= \E\left[\sum_{t=1}^T \left( \KL(\delta_{\phi^\star}, \rho_t) - \KL(\delta_{\phi^\star},  (\nu_t)_{\bo{\phi}}(\cdot|\pi_t,o_t))\right) \right] = \E\left[ \sum_{t=1}^T \log\frac{\nu_t(\phi^\star|\pi_t,o_t)}{\rho_t(\phi^\star)} \right],   
\end{align*}
where in the second equality we use that $o_t$ is drawn from $M_t(\cdot|\pi_t)$. 
%Specifically, if we choose $F_1(\pi, \nu, \rho_t) = \frac{1}{\eta}\E_{(M,\pi^\star) \sim \nu}\E_{z \sim M(\pi)}\left[\KL(\nu_{\bo{\phi}}(\cdot|\pi, z), \rho_t)\right]$ and $F_2(\rho; \pi_t, z_t) = \KL\left(\rho, \nu_{\bo{\phi}}(\cdot|\pi_t, z_t)\right)$, then the $\air^{\Phi}_{F_1, \rho_t}(p,\nu)$ is exactly the standard $\Phi$-AIR defined in \pref{eq:phiair}. Moreover, $\rho_{t+1} = \argmin_{\rho \in \Delta(\Phi)} F_2(\rho, \pi_t, z_t) = \nu_{\bo{\phi}}\left(\cdot|\pi_t, z_t\right)$. This exactly recovers the algorithm in \cite{liu2025decision}. Since the Bregman divergence $D_{F_1}(\nu, \nu') = \E_{(M,\pi^\star) \sim \nu}\E_{z \sim M(\pi)}\left[\KL\left(\nu_{\bo{\phi}}(\cdot|\pi, z), \nu'_{\bo{\phi}}(\cdot|\pi, z)\right)\right]$. Let $\phi_t^\star$ be the aggregation that $(M_t, \pi^\star)$ lies in,  the overhead in \pref{eq:minimax-guarantee} now becomes 
%\begin{align}
%    \frac{1}{\eta}\E_{\pi \sim p_t}\E_{o \sim M_t(\pi)}\left[\log\left(\frac{\nu_t(\phi^\star_t|\pi, z)}{\rho_t(\phi_t^\star)}\right)\right] =  \frac{1}{\eta}\E_{\pi \sim p_t}\E_{z \sim M_t(\pi)}\left[\log\left(\frac{\nu_t(\phi^\star_t|\pi, z)}{\nu_{t}(\phi_t^\star|\pi_t, z_t)}\right)\right] +  \frac{1}{\eta}\log\left(\frac{\rho_{t+1}(\phi_t^\star)}{\rho_t(\phi_t^\star)}\right) \label{eq:special-overhead}
%\end{align}
Thus, by letting $\rho_{t+1}(\phi) = \nu_t(\phi|\pi_t,o_t)$, their algorithm achieves $\E[\est] = \E\left[\sum_{t=1}^T \log\frac{\rho_{t+1}(\phi^\star)}{\rho_t(\phi^\star)} \right]\leq \log\frac{1}{\rho_1(\phi^\star)}=\log|\Phi|$. 
%For all the learnable instances discussed in \cite{xu2023bayesian} and \cite{liu2025decision}, $\phi_t^\star = \phi^\star$ for every $t \in [T]$ where $\phi^\star$ is a fixed aggregation. For instance, following the discussion on $\Phi$ in \pref{sec:pre},  model-based RL in stochastic settings has $M_t = M^\star$ and $\pi^\star = \pi_{M^\star}$, thus $\phi^\star = \phi_{M^\star}$, while model-free RL has  $\phi^\star = \phi_{Q^\star}$ where $Q^\star$ is the optimal $Q$ function of model $M^\star$.  For adversarial bandit, $M_t$ is changing over time but $\pi^\star$ is fixed, so $\phi^\star = \phi_{\pi^\star}$. For hybrid MDPs, transition $P^\star$ and comparator policy $\pi^\star$ are fixed over episodes so $\phi^\star = \phi^\star_{P^\star, \pi^\star}$. Given $\phi_t^\star = \phi^\star$, putting 
Using this in \pref{eq:minimax-guarantee} proves \pref{thm: air phi}. The results of \cite{xu2023bayesian} can also be recovered as they are special cases of \cite{liu2025decision}. 
%\begin{align*}
%    \E[\Reg] \le T \cdot \max_{\rho \in \Delta(\Phi)}\min_{p \in \Delta(\Pi)}\max_{\nu \in \Delta(\Psi)}\air^{\Phi}_{\rho, \eta}(p,\nu) + \frac{1}{\eta}\log(|\Phi|)
%\end{align*}
%which is exactly the regret bound derived in \cite{liu2025decision} and cover the results in \cite{xu2023bayesian}. The complexity measure based on $\Phi$-AIR is tightly related with DEC and uncover the learnability of numerous instances in stochastic, adversarial and hybrid settings. However, the $\Phi$-AIR approaches still fall short to deal with hybrid settings with full information feedback and model-free learning in either stochastic or hybrid settings. Our framework \pref{alg:general} can tackle both of these challenges with carefully designed $F_1$ and $F_2$.

\subsection{Recovering results for adversarial MDP with full-information feedback \citep{liu2025decision}}  
For learning with full information feedback in the adversarial MDPs, the learner can observe the full reward function at the end of each episode. In other words, at episode $t$, the reward function $R_t: \calS\times \calA\to [0,1]$ is part of the observation $o_t$. In this setting, the $\log|\Pi|$ dependence in the regret bound can be improved to $\log|\calA|$. To achieve this, \cite{liu2025decision} designed a two-level algorithm and define a new notion called InfoAIR. We can recover this result by instantiating our \pref{alg:general} with $\Phi=\{\phi_{P,(a_s)_{s\in\calS}}:~ P\in\calP, a_s \in\calA, \forall s\in\calS\}$ where $\phi_{P,(a_s)_{s\in\calS}} = \{((P, R), \pi^\star): R\in\calR, \pi^\star=(a_s)_{s\in\calS}\}$, that is, partitioning $\calM\times \Pi$ according to the transition kernel and the actions taken by the policy on all states.  
Then define 
\begin{align*}
   D^\pi(\nu\| \rho) = \E_{(P,R,\pi^\star) \sim \nu}\E_{o \sim M_{P,R}(\cdot|\pi)}\E_{s \sim d^{\pi, P}}\left[\KL(\nu_{\bo{a}_s, \bo{P}}(\cdot|\pi, o), \rho_{\bo{a}_s, \bo{P}})\right], 
\end{align*}
where $M_{P, R}$ denotes the MDP model with transition kernel $P$ and reward function $R$, and $\rho_{\bo{a}_s, \bo{P}}$ denotes $\rho$'s marginal distribution over $(a_s, P)$ following our notational convention. Finally, update the posterior as 
%where $\nu(a, P|\pi, o, s) = \sum_{\pi^\star \in \Pi}\pi^\star(a|s)\nu(\pi^\star, P|\pi, o, R, s)$, and  
$\rho_{t+1} = \argmin_{\rho}\sum_{s \in \calS}\KL\left(\rho_{\bo{a}_s, \bo{P}}, \nu_{\bo{a}_s, \bo{P}}(\cdot|\pi_t, o_t)\right)$. This recovers the same regret bound as in \cite{liu2025decision} without the need for the two-level design. %Such design of $F_1$ and $F_2$ take advantages of the full information structure by putting the whole reward function into the condition of $\nu$. This enables the belief $\nu$ can be updated in the action space $\calA$ at every state instead of the whole policy space $\Pi$, which is analog to policy optimization approaches \citep{sherman2023rate, liu2023towards, casselwarm}. 
We also note that the analysis for this result requires our new proof strategy in \pref{eq:general-decom}, as the $D^\pi(\nu\|\rho)$ here is not strictly convex in $\nu$ and the previous proof \cite{xu2023bayesian, liu2025decision} cannot be applied. %\cw{The notation in this paragraph is confusing ($\rho(\cdot|s)?$) and there is no definition of $\Phi$ here. }




%More importantly, for model-free learning, standard $\Phi$-AIR can only handle very restricted instance in stochastic settings and fails to be extended to hybrid settings. \cite{liu2025decision} design a separate algorithm based on optimistic DEC to handle model-free learning in hybrid settings with full information feedback. Their approach, however, cannot be extended to bandit feedback settings due to technical limitations. We will show that, with carefully chosen of $F_1, F_2$ in our framework \pref{alg:general}, we not only improve the standard optimistic DEC designed for model-free RL in stochastic settings (\pref{sec:stochastic}), but also, for the first time, provide an approach to deal with model-free learning in hybrid settings with general function approximation (\pref{sec:hybrid}).  


